% Architecture and experimental-setup section.
% Intended to be \input into the main paper (methodology.tex).
% The results table is generated by paper_data/architecture_table.py from the
% raw metrics in results_before_IMPALA/ — rerun that script to audit or
% refresh every number reported here.

\section{Architecture and Training Setup}
\label{sec:architecture}

All agents share a single recurrent Q-network: a convolutional encoder, an
LSTM core, and one linear head that carries both the value estimate and the
policy readout. The imitation loss applied to this head is described in
Section~\ref{sec:methodology} and is excluded here.

\subsection{Convolutional Encoders}

\paragraph{Nature CNN.} The baseline encoder is the three-layer network of
Mnih et al.\ (2015): $32$ filters of $8{\times}8$ stride~4, $64$ filters of
$4{\times}4$ stride~2, and $64$ filters of $3{\times}3$ stride~1, followed by
a fully connected projection to a $512$-dimensional feature vector. All
activations are ELU rather than ReLU; with backpropagation through time,
ReLU units in the encoder died frequently enough to destabilize early
training. The encoder has ${\sim}0.60$M parameters. Every result in
Table~\ref{tab:nature-results} uses this encoder.

\paragraph{IMPALA ResNet.} Later experiments (not reported in
Table~\ref{tab:nature-results}) replace the encoder with the residual
network of Espeholt et al.\ (2018): three stages of depth $16$, $32$, and
$32$, where each stage applies a $3{\times}3$ convolution, a $3{\times}3$
max-pool of stride~2, and two pre-activation residual blocks (each
ReLU--conv--ReLU--conv with an additive skip). A $64{\times}64$ observation
leaves the stages at $8{\times}8{\times}32$ and is projected by a ReLU--FC
layer to the same $512$-dimensional feature space (${\sim}1.15$M
parameters), so the recurrent core is identical under either encoder.

\subsection{Recurrent Core and Dueling Head}

Encoder features feed a single-layer LSTM with $512$ hidden units. A single
linear layer maps the LSTM output to $|\mathcal{A}|{+}1$ values, decomposed
dueling-style into a state value $V$ and advantages $A$ with
$Q = V + (A - \bar{A})$. The same advantage stream doubles as the policy
logits for the imitation term, so value learning and behavior cloning share
every parameter of the network.

\subsection{Recurrent Value Learning}

Training follows R2D2 (Kapturowski et al., 2019). Replay stores whole
episodes in flat, pre-allocated GPU buffers ($150$k online and $50$k
demonstration transitions, observations kept as \texttt{uint8}); each
transition also stores the actor's LSTM state at collection time in
\texttt{fp16}. Updates sample windows of $64$ steps, re-initialize the LSTM
from the stored state, and burn in the first $16$ steps without gradient.
Demonstration episodes carry no actor states, so their stored states are
refreshed by replaying each episode through the current network once per
iteration. Targets use double DQN (online-network argmax, target-network
evaluation) with a soft target update ($\tau = 0.005$) and discount
$\gamma = 0.99$; where noted, the demonstration channel uses uncorrected
$5$-step returns while online data remains $1$-step. Optimization is Adam
(learning rate $3{\times}10^{-4}$) with a global gradient-norm clip of
$1.0$. Exploration is $\epsilon$-greedy, annealed $0.25 \rightarrow 0.05$
over the first $50$k frames; evaluation acts greedily.

Experience is collected from $8$ Crafter environments running in parallel
worker processes, $2{,}000$ frames per iteration, with $30$ update epochs
(batches of $64$ windows) per iteration. All runs in
Table~\ref{tab:nature-results} train for $400$ iterations ($1.2$M
environment frames). Later IMPALA-encoder experiments additionally use
prioritized experience replay on the online buffer ($\alpha{=}0.6$,
$\beta{=}0.4$, priorities taken from the TD errors of each update) and
combine all loss terms into a single backward pass per epoch.

\subsection{Demonstration Data and Integration Mechanisms}

The human dataset contains $51$ episodes ($37{,}403$ frames) recorded by a
single player, mean return $18.2$. The methods in
Table~\ref{tab:nature-results} differ only in how this data enters
training:

\begin{itemize}
  \item \textbf{Behavior cloning / RCQL (offline):} supervised imitation or
  conservative Q-learning (Kumar et al., 2020) on the demonstrations alone,
  with no environment interaction.
  \item \textbf{R2D2 (online only):} no demonstration data.
  \item \textbf{R2D3:} demonstrations enter as plain TD replay on a
  $1/16$ demo-to-online ratio (Paine et al., 2019) with $5$-step returns.
  \item \textbf{CQL anchor (annealed):} every epoch trains a CQL loss on a
  demonstration batch alongside online TD; the CQL weight anneals
  $1 \rightarrow 0$ over $500$k frames.
  \item \textbf{BC term (annealed):} the imitation loss of
  Section~\ref{sec:methodology} on a demonstration batch every epoch,
  weight annealed $1 \rightarrow 0.1$ over $500$k frames.
  \item \textbf{DQfD-lite:} the annealed BC term plus the R2D3 demo TD
  channel ($1/16$, $5$-step), in the spirit of DQfD (Hester et al., 2018).
  \item \textbf{Demo starts:} $2$ of the $8$ collection environments reset
  from mid-demonstration states instead of fresh worlds (cf.\ Backplay,
  Resnick et al., 2018). Crafter states are reconstructed exactly from the
  recorded semantic world map, inventory, achievement set, and player
  position; only achievements the demonstrator had not yet unlocked yield
  reward. Restart points are sampled with priority toward states just
  before large TD errors on the demonstration, re-scored each iteration.
\end{itemize}

\subsection{Evaluation}

Agents are evaluated greedily on $5$ freshly generated worlds per
iteration. Table~\ref{tab:nature-results} pools evaluation episodes from
iterations $301$--$400$ (the matched $1.2$M-frame budget; $\geq 495$
episodes per method) and reports the mean episode return and the official
Crafter score (Hafner, 2021): with $s_i$ the success percentage of
achievement $i$ over the pooled episodes,
$S = \exp\!\big(\tfrac{1}{22}\sum_{i=1}^{22}\ln(1+s_i)\big) - 1$. The
geometric mean rewards broad progress through the achievement hierarchy
rather than repetition of easy achievements, and is the statistic under
which the demonstration-integration methods separate most clearly. Offline
baselines are evaluated at the end of their (interaction-free) training;
the human row scores the $51$ demonstration episodes themselves. The
online-only R2D2 runs predate achievement logging, so their Crafter score
is unavailable.

\begin{table}[h]
\centering
\input{paper_data/architecture_table}
\caption{Crafter results with the Nature CNN encoder at the matched
$1.2$M-frame budget (evaluation episodes pooled from iterations
$301$--$400$; mean $\pm$ std of per-evaluation returns). Offline baselines
use no environment interaction. Generated by
\texttt{paper\_data/architecture\_table.py}. Published reference scores at
$1$M frames: PPO $4.6\%$, Rainbow $4.3\%$, DreamerV2 ${\sim}10\%$
(Hafner, 2021).}
\label{tab:nature-results}
\end{table}
